% 正方晶格实空间与倒空间示意；信息按“几何—路径—对偶关系”分层。
\resizebox{\textwidth}{!}{%
\begin{tikzpicture}[x=1cm,y=1cm,>=Latex,font=\small]
  \tikzset{
    panel/.style={draw=gray!45,rounded corners=3pt,fill=gray!2},
    pointlabel/.style={font=\footnotesize,fill=white,inner sep=1.5pt}
  }

  \begin{scope}[shift={(-4.2,0.35)}]
    \draw[panel] (-3.0,-2.65) rectangle (3.0,2.75);
    \node[font=\bfseries,anchor=west] at (-2.75,2.4) {(a) 实空间正方晶格};
    \foreach \x in {-2,-1,0,1,2} {
      \foreach \y in {-2,-1,0,1,2} {\fill (\x,\y) circle (2pt);}
    }
    \draw[green!60!black,very thick] (-0.5,-0.5) rectangle (0.5,0.5);
    \node[green!45!black,anchor=north] at (0,-2.95) {Wigner--Seitz 原胞};
    \draw[->,very thick,blue] (0,0) -- (1.45,0)
      node[at end,below=8pt] {$\mathbf a_1=a\hat{\mathbf x}$};
    \draw[->,very thick,red] (0,0) -- (0,1.45)
      node[midway,left=3pt] {$\mathbf a_2=a\hat{\mathbf y}$};
    \draw[<->,thick] (2.35,-0.5) -- (2.35,0.5) node[midway,right] {$a$};
  \end{scope}

  \begin{scope}[shift={(4.15,0.35)}]
    \draw[panel] (-3.35,-2.65) rectangle (3.35,2.75);
    \node[font=\bfseries,anchor=west] at (-3.1,2.4) {(b) 第一布里渊区与高对称路径};
    \fill[blue!9] (-1.45,-1.45) rectangle (1.45,1.45);
    \draw[very thick,blue!70] (-1.45,-1.45) rectangle (1.45,1.45);
    \draw[->,thick] (-2.05,0) -- (2.65,0) node[right] {$k_x$};
    \draw[->,thick] (0,-2.0) -- (0,1.65) node[left] {$k_y$};

    \fill (0,0) circle (2.2pt);
    \fill (1.45,0) circle (2.2pt);
    \fill (1.45,1.45) circle (2.2pt);
    \node[pointlabel,anchor=north east] at (-0.05,-0.06) {$\Gamma(0,0)$};
    \node[pointlabel,anchor=north west] at (1.60,-0.10) {$X$};
    \node[pointlabel,anchor=south west] at (1.55,1.52) {$M$};

    \draw[very thick,red!75!black,->] (0.12,0) -- (1.28,0);
    \draw[very thick,red!75!black,->] (1.45,0.14) -- (1.45,1.27);
    \draw[very thick,red!75!black,->] (1.34,1.34) -- (0.12,0.12);
    \node[red!70!black,font=\footnotesize,anchor=north] at (.9,-.65) {$\Gamma\to X$};
    \node[red!70!black,font=\footnotesize,anchor=west] at (1.68,0.72) {$X\to M$};
    \node[red!70!black,font=\footnotesize,fill=white,inner sep=1pt] at (-.60,1.15) {$M\to\Gamma$};

    \node[font=\footnotesize,anchor=south] at (0,-2.35)
      {常用能带路径：$\Gamma\rightarrow X\rightarrow M\rightarrow\Gamma$};
  \end{scope}

  \draw[->,very thick,purple!70!black] (-1.1,3.35) -- (1.1,3.35)
    node[midway,above] {傅里叶变换};
  \node[draw=gray!55,rounded corners,fill=white,align=center] at (0,-3.8)
    {$\mathbf b_i\!\cdot\!\mathbf a_j=2\pi\delta_{ij}$\qquad
     $\mathbf b_1=\dfrac{2\pi}{a}\hat{\mathbf x}$\qquad
     $\mathbf b_2=\dfrac{2\pi}{a}\hat{\mathbf y}$};
\end{tikzpicture}%
}
